Throughout, walls are tame (\claimref{sheaf.tameness}), as in the results cited.

\emph{(i).} By clause (1) of \claimref{synthesis.definite-side-rung-count-on-the-quarter-turn-lens-wall-of-dimension-five}, on $\Sig_{n}$ the number $e(\shL)$ is the least $k$ with a continuous $F : S^{2n+1} \to \R^{k} \setminus \{0\}$, $F(J(v)) = -F(v)$, with $J$ the quarter turn in every coordinate plane of $\R^{2n+2}$. This is $s(m)$ of \claimref{synthesis.imported-stolz-level-of-real-projective-spaces-under-the-quarter-turn} with $2m - 1 = 2n + 1$, $m = n + 1 \ge 2$. The residues: $m \equiv 2, 0$ exactly when $n \equiv 1, 7$, and then $s = m + 1 = n + 2$; $m \equiv 6$ exactly when $n \equiv 5$, and then $s = m + 3 = n + 4$; otherwise $s = m + 2 = n + 3$. On a definite side $k_{+} = e(\shL)$ by the same cited result (its opening paragraph, from clauses (1) and (3) of the two-rung result it cites). For $n = 1, \dots, 8$: $m = 2, \dots, 9$ and $s = 3, 5, 6, 7, 9, 9, 9, 11$. The cited result has $k_{+}(\Sig_{1}) \le 3$ with the stored value $3$ for $L(4,1)$, $k_{+}(\Sig_{2}) = 5$ and $5 \le k_{+}(\Sig_{3}) \le 6$.

\emph{(ii).} A completion with $k$ rungs exists when $\shL$ sits in a trivial bundle of rank $k$ (clause (2) of \claimref{synthesis.wall-rung-count-is-height-of-dying-line-class}), so $k(\Sig) \le e(\shL)$ on every side, and $e(\shL) = s(n+1)$ by (i). The lower bound is clause (2) of \claimref{synthesis.lens-space-walls-need-more-rungs-than-the-mod-two-height} at its version 3, which holds given the lambda-operations. The bound $k \ge 4$ for $n \ge 2$ is clause (3) of \claimref{synthesis.three-rungs-do-not-complete-the-quarter-turn-lens-walls-on-any-side}, conditional at $n = 2$ on the spin class as stated there; with $s(3) = 5$ and $s(4) = 6$ it gives the two windows. Since $s(n+1) \ge n + 2$ and $1 + \lceil n/2 \rceil \le (n + 3)/2$, the ratio of the two tends to $\tfrac12$.

\emph{(iii).} With $m = n + 1$: $\lceil 3m/2 \rceil - m = \lceil m/2 \rceil$, and $s(m) - m \in \{1, 2, 3\}$, with the value $3$ only for $m \equiv 6$ modulo $8$. The excess of the explicit bound over $s(m)$ is $\lceil m/2 \rceil - (s(m) - m)$. For $m \ge 7$ with $m \not\equiv 6$: $\lceil m/2 \rceil \ge 4$ and $s(m) - m \le 2$, so the excess is at least $2$. For $m \ge 7$ with $m \equiv 6$: $m \ge 14$, $\lceil m/2 \rceil \ge 7$, and the excess is at least $4$. For $m = 2, \dots, 6$ the pairs $(\lceil 3m/2 \rceil, s(m))$ are $(3,3)$, $(5,5)$, $(6,6)$, $(8,7)$, $(9,9)$: equality at $m = 2, 3, 4, 6$, that is $n = 1, 2, 3, 5$, and excess $1$ at $n = 4$.

\emph{(iv).} Read off (i): $k_{+} - n$ is a function of $n + 1$ modulo $8$. The period of the real Clifford algebras is \claimref{std-algebra.periodicity}. No relation between the two is proved or used.
